\documentclass[11pt]{article}
\usepackage[T1]{fontenc}
\usepackage{lmodern,amsmath,amssymb,amsthm,microtype}
\usepackage[margin=1in]{geometry}
\usepackage[hidelinks]{hyperref}
\newtheorem{theorem}{Theorem}
\newtheorem{lemma}[theorem]{Lemma}
\newcommand{\R}{\mathbb R}
\newcommand{\Z}{\mathbb Z}
\newcommand{\Q}{\mathbb Q}
\newcommand{\conv}{\operatorname{conv}}
\title{Polynomial-time optimization and exact separation\newline
for rational mixed-integer bipartite covers}
\author{Alec Kriebel\thanks{Independent researcher. ORCID:
\href{https://orcid.org/0009-0001-9320-500X}{0009-0001-9320-500X}.}}
\date{12 October 2026\quad (research note)}
\begin{document}
\maketitle
\begin{abstract}
We give polynomial-time exact linear optimization, strong separation and
membership for mixed-integer bipartite covers with arbitrary rational edge
demands, unrestricted variable signs and any selected integer coordinates.
All data and query points are rational in the total binary-input model;
dimension and demand denominators vary. This resolves the finite rational
interpretation of Conforti's 2008 hull-membership problem. A direct specialized
proximity proof localizes integer labels to intervals of length $2n$ and
proves attainment. Exact continuous completion yields a submodular value
oracle on a polynomial-size threshold lattice. A bounded-core argument
establishes closedness of the actual hull and a polynomial facet-encoding
bound, allowing classical optimization--separation equivalence to apply.
The result is an explicit application of classical mixed proximity,
discrete-convex and submodular-minimization ideas. It does not prove the
separate subtree-hull conjecture or assert a compact formulation or strongly
polynomial total algorithm.
\end{abstract}

\section{Problem and contribution}
Let $U,V$ be specified disjoint vertex classes of a finite bipartite graph,
$n=|U|+|V|$, $E\subset U\times V$, $I\subset U\cup V$, and $b\in\Q^E$.
Define
\[
 S=\{x\in\R^n:x_u+x_v\ge b_{uv}\ (uv\in E),\ x_i\in\Z\ (i\in I)\}.
\]
There is no nonnegativity assumption. Numerators, positive denominators,
graph and integer-coordinate indices are encoded explicitly in binary.
Conforti \cite[pp.~2904--2905, Problem~1]{OWR} asks the complexity of membership
in $\conv(S)$ and discusses linear optimization. Since that report leaves
encoding of real query points implicit, our theorem specifies rational query
points in the standard finite bit model.

\begin{theorem}\label{thm:main}
For the data above and rational objective $c$, in polynomial total binary-input
time one can return an exact rational minimizer of $c^Tx$ on $S$, or a feasible
point and an integral improving recession direction if the infimum is
$-\infty$. The actual hull $\conv(S)$ is a closed rational polyhedron with a
polynomially computable facet-encoding bound. Strong separation and exact
membership at every rational query, including boundary queries, are
polynomial-time tasks. All $n$, $|I|$ and denominators may vary.
\end{theorem}

The zero-dimensional case is immediate; assume $n\ge1$ below. Empty edge or
integer-coordinate sets are allowed. Feasibility follows by setting every
coordinate to a common integer at least zero and all demands. If $I$ is empty,
ordinary rational LP gives optimization directly; the hull argument below
still applies.

Mixed proximity is classical: Cook--Gerards--Schrijver--Tardos
\cite[p.~258, Remarks~(1)]{CGST} give all-coordinate distance at most $n\Delta$
for integral combined constraint matrices, assuming a mixed optimum exists.
Our signed-threshold proof is a direct specialization, and supplies attainment
without that premise. Hybrid partial minimization and discrete-convex
optimization are prior mechanisms \cite[Proposition~2]{THM};
\cite[\S4.1]{Murota} documents an earlier $n^2$-element ring-family lineage.
We claim no new general proximity or abstract ring-minimization method.
The contribution is the explicit rational application, exact LP completion
oracle, denominator-independent label count and actual-hull bit accounting.

Earlier network and bipartite-cover formulations \cite{Network,Covers}
depend polynomially on a numerical common denominator and explicitly discuss
the remaining binary-size gap. They do not justify enumerating denominator
residues here. A bounded source audit located no earlier full solution to
this exact claim in inspected bodies; exhaustive priority is unverified.

\section{Difference variables and proximity}
Put $y_u=x_u$ for $u\in U$, $y_v=-x_v$ for $v\in V$, and transform $c$ to $a$
by the same sign change. This integral invertible map preserves the selected
integer coordinates. Write
\[
 P=\{y:y_u-y_v\ge b_{uv}\},\qquad
 T=P\cap\{y:y_i\in\Z\ (i\in I)\}.
\]
Rational LP decides boundedness. If $a^Ty$ is unbounded on $P$, rational
recession theory supplies a polynomial-bit rational solution of
\[
 r_u-r_v\ge0\quad(uv\in E),\qquad a^Tr\le-1.
\]
Clearing coordinate denominators produces a polynomial-bit integral improving
ray. From the integral feasible point, every nonnegative integer step stays
in $T$. Thus boundedness on $P$ and $T$ agrees. Otherwise rational LP returns
a polynomial-bit minimizer $y^*$ on $P$. No vertex of this unanchored,
generally nonpointed polyhedron is assumed.

\begin{lemma}\label{lem:proximity}
If $a$ is bounded below on $P$, the minimum on $T$ is attained, and at least
one minimizer $z$ satisfies $\|z-y^*\|_\infty<n$.
\end{lemma}
\begin{proof}
Every $d\in\R^n$ has a signed nested-threshold decomposition
$d=\sum_{h=1}^s\lambda_h r^h$, with $s\le n$, $\lambda_h>0$ and nonzero
$r^h\in\{0,1\}^n$ or $\{0,-1\}^n$. For positive coordinates, list distinct
values $0<t_1<\cdots<t_j$, use indicators of $\{i:d_i\ge t_h\}$ and successive
differences $t_h-t_{h-1}$; do the same with negative magnitudes and negative
indicators. Nesting makes every layer conformal on coordinates and differences:
$r_i^h$ has the sign of $d_i$ or is zero, and $r_i^h-r_j^h$ has the sign of
$d_i-d_j$ or is zero. Consequently each weighted coordinate and difference
lies between zero and the corresponding value of $d$.

Take $z\in T$ and $d=z-y^*$. For each layer $(\lambda,r)$, the point
$y^*+\lambda r$ lies in $P$, rowwise between $y^*$ and $z$. Optimality gives
$a^Tr\ge0$. If $\lambda\ge1$, the point $z-r$ also lies rowwise between those
endpoints, stays mixed-integer, has no greater objective and has strictly
smaller $\ell_1$ distance to $y^*$.

To avoid assuming attainment, fix any $z^0\in T$. The closed nonempty sublevel
set $A=\{z\in T:a^Tz\le a^Tz^0\}$ has a distance minimizer: intersect it with
the closed $\ell_1$ ball of radius $\|z^0-y^*\|_1$. Compactness and the removal
step rule out every layer of weight at least one at this minimizer. With at
most $n$ layers, its sup-distance is less than $n$.
Hence every point of $T$ has a no-more-expensive representative in the same
compact box $y^*+[-n,n]^n$. Minimization on its closed mixed subset attains
the global optimum, and a distance-minimizing optimum has the asserted strict
bound.
\end{proof}

\section{Exact partial minimization on a short-label lattice}
For $i\in I$ set
\[
 l_i=\lceil y_i^*-n\rceil,\quad h_i=\lfloor y_i^*+n\rfloor,\quad
 D=\{p\in\Z^I:l_i\le p_i\le h_i,\
 p_u-p_v\ge\lceil b_{uv}\rceil\text{ if }u,v\in I\}.
\]
Each width is at most $2n$. Exactly these assignments have completions in $P$:
necessity is immediate; for sufficiency undo the sign change on fixed
coordinates and set every remaining original coordinate to a sufficiently
large common integer. Only edges between two integer coordinates are rounded.
The free-sign bipartite model makes this completion possible. By
Lemma~\ref{lem:proximity}, $D$ is nonempty.

For $p\in D$, the slice value
\[
 f(p)=\min\{a^Ty:y\in P,\ y_i=p_i\ (i\in I)\}
\]
is finite, attained and computable exactly by polynomial-size rational LP.
Both $P$ and $D$ are closed under coordinatewise meet and join. For example,
two inequalities $y_u\ge y_v+b$, $z_u\ge z_v+b$ imply their min and max
counterparts. Optimal completions $y,z$ of $p,q$ give
\[
 f(p)+f(q)=a^Ty+a^Tz
 =a^T(y\wedge z)+a^T(y\vee z)
 \ge f(p\wedge q)+f(p\vee q).
\]
Thus $f$ is submodular on the finite distributive lattice $D$. This closure
is in difference variables; it would be false for original cover variables.

\paragraph{Uniform rational encoding.}
Let $B=\max(1,\max_e\lceil|b_e|\rceil)$,
$A_0=\max(1,\max_i\lceil|a_i|\rceil)$,
$H^*=\lceil\|y^*\|_\infty\rceil+n$, and let $k,q$ be the lcms of demand and
objective denominators (empty lcm is one). Every $p$ has magnitude at most
$H^*$. Anchor one coordinate at zero in each slice component with no fixed
coordinate; boundedness makes its objective invariant under translation.
The resulting slice is pointed. At an optimal vertex each tight-edge
component connects to a fixed or zero anchor, since otherwise it could be
translated locally, contradicting full active rank. A tight spanning forest
therefore expresses every coordinate as an anchor plus at most $n-1$ signed
demands. Its magnitude is at most $H^*+nB$ and its denominator divides $k$.
Consequently
\[
 |f(p)|\le nA_0(H^*+nB),\qquad \operatorname{den}(f(p))\mid kq.
\]
The value encoding bound is uniform and polynomial in binary input size;
$\log k$ and $\log q$ are at most sums of input denominator lengths.

\section{Polynomial-size domain interface and minimization}
Make a node $(i,t)$ for each $i\in I$, $l_i<t\le h_i$, meaning $p_i\ge t$.
The ground set $W$ has size at most $2n|I|$. Encode
$A(p)=\{(i,t):l_i<t\le p_i\}$. Union/intersection correspond to join/meet.
The feasible family $\mathcal C=\{A(p):p\in D\}$ has explicit implications:
\begin{enumerate}
\item Prefix arcs $(i,t)\to(i,t-1)$ when $t>l_i+1$.
\item For an integer--integer edge put $d=\lceil b_{uv}\rceil$. The base
condition $p_u\ge l_v+d$ is automatic if $l_v+d\le l_u$, forces node
$(u,l_v+d)$ if in range, and makes the family empty if above $h_u$.
\item Each node $(v,t)$ implies $(u,t+d)$. Targets at most $l_u$ are automatic;
targets above $h_u$ forbid $(v,t)$; other targets give arcs.
\end{enumerate}
These conditions are equivalent to $p\in D$, including the essential base
condition when $p_v=l_v$. There are $O(n^2+n|E|)$ arcs/nodes and polynomial-bit
labels, independently of numerical denominators.

A member contains all forced nodes, avoids all forbidden nodes and is closed
under arcs. Reachability computes the smallest member and the smallest member
containing any given node, or detects its nonexistence. This is the ring-family
interface of Schrijver \cite[\S6]{Schrijver}, supplied explicitly rather than
assumed from a value oracle. Equivalently compute the intersection and union
of members and the closure digraph on the remaining nodes, giving the lattice
encoding of \cite[\S10.3]{GLS}.

Apply established exact rational submodular minimization to
$F(A(p))=f(p)$. The uniform encoding bound above and
\cite[Theorem~10.2.7, problem~10.3.8 and its reduction]{GLS} give polynomial
oracle-bit complexity on this lattice; \cite[\S6]{Schrijver} supplies the
classical combinatorial ring reduction as well. Counting arithmetic operations
alone would not suffice. All oracle calls here are polynomial-bit rational LPs.
A minimizing assignment and one optimizing completion give a global mixed
minimizer by Lemma~\ref{lem:proximity}. The full algorithm is polynomial in
binary length; a strongly polynomial total LP-based algorithm is not claimed.

\section{The actual closed hull and facet complexity}
Let $H=nB+n$, $R=\{r:r_u-r_v\ge0\}$ and
$K=\conv(T\cap[-H,H]^n)$. These are mathematical objects for a size argument,
not an enumeration performed by the algorithm.

For every objective bounded below on $P$, the objective annihilates component
translation lineality. Anchor one coordinate per graph component at zero.
An optimal vertex of this pointed relaxation has tight spanning trees
connecting each component to its anchor. Its coordinates are sums of at most
$n-1$ signed demands and have norm at most $nB$. Lemma~\ref{lem:proximity}
puts a mixed optimum within norm $H$, uniformly over all bounded objectives.
Thus every such objective attains its minimum on $T$ inside $K$.

The bounded mixed set is a finite union of bounded rational slice polytopes,
so $K$ is a rational polytope and $K+R$ is a closed rational polyhedron.
Every $r\in R$ decomposes into nonnegative multiples of integral signed
threshold layers in $R$. For such an integral layer $d$ and $z\in T$,
$z+md\in T$ for all nonnegative integers $m$. Interpolation of consecutive
steps puts $z+td$ in $\conv(T)$ for every real $t\ge0$. The same holds for
finite convex combinations, then finite sums of layers. Hence $K+R\subset
\conv(T)$.

If some $z\in T$ lay outside $K+R$, strict separation would give $a$ with
$a^Tz<\inf_{w\in K+R}a^Tw$. The right side is finite, so $a$ is nonnegative
on $R$ and bounded below on $P$. The uniform small-optimum conclusion makes
its minimum on $T$ equal its minimum on $K$, a contradiction. Therefore
\[
 \conv(T)=K+R,
\]
which proves closedness of the actual hull, without substituting its closure.

Slice matrices consist of directed incidence rows, unit fixings and integer
box bounds, and are totally unimodular. Each slice vertex lies in
$(1/k)\Z^n$ and has norm at most $H$. Thus $K$ is generated by finitely many
such vertices; $R$ is generated by integral threshold rays of norm at most
one. The hull is full-dimensional, since original cover variables contain a
translated positive orthant in the convex hull of large integral vectors.
Homogenize point and ray generators as $(v,1),(r,0)$ and multiply point
generators by $k$. Integer generator coordinates have magnitude at most
$k(H+1)$. Every facet has $n$ independent homogeneous generators on it;
cofactors give an integral facet inequality, including constant term, with
coefficients bounded by $n![k(H+1)]^n$. Lineality causes no problem: spanning
generators may include both signs of lineality directions.

A computable generous bound on the total signed binary encoding of each
facet inequality is
\[
 \varphi=10(n+1)^2\bigl[1+\lceil\log_2(n+1)\rceil+
 \lceil\log_2(k+1)\rceil+\lceil\log_2(H+2)\rceil\bigr].
\]
It dominates the cofactor coefficient lengths and is polynomial in input
encoding. No vertices, facets or denominator residues are enumerated.

\section{Strong separation, scope and validation}
The optimization procedure returns exact feasible optimizing points or
improving integral recession directions. Points need not be vertices.
Rescale an improving ray to the normalization in the strong-optimization
interface of \cite[Definitions~6.2.1--6.2.2]{GLS}. Together with $\varphi$,
this gives a well-described rational polyhedron. The exact strong
optimization--separation equivalence \cite[Theorem~6.4.9, pp.~179--180]{GLS}
therefore provides a polynomial-time strict separator or exact membership
confirmation at any rational query. Undo the diagonal sign change to obtain
Theorem~\ref{thm:main}. Unbounded polyhedra and boundary queries are included;
no approximate membership or closure surrogate is used.

The separate subtree-hull Conjecture~2 of \cite{OWR} remains unresolved here.
No polynomial-size extended formulation, facet classification, arbitrary-real
oracle result or full polynomial-optimizer implementation is claimed.
The finite exact checkers enumerate only small cases; they do not implement
LP, submodular minimization and GLS in polynomial time or prove complexity
from observed runtime. Independent adversarial families checked proximity,
attainment and oracle bits, and ring/hull/separation interfaces. The supplement
records primary-source comparisons, exact diagnostics and limitations of the
bounded priority audit. Classical ingredients and earlier formulations receive
credit above; no first-ever or exhaustive novelty certificate is asserted.

AI tools were used extensively in solving, verifying, source auditing and
preparing this note and supporting materials. This is an unrefereed preprint.
Independent adversarial reviews are AI reviews; conventional human peer review,
journal acceptance and formal proof-assistant certification are not claimed.

\begin{thebibliography}{9}
\bibitem{OWR} M.~Conforti, \emph{Combinatorial Mixed-Integer Programming},
Oberwolfach Reports 51/2008, 2904--2905.
\href{https://ems.press/content/serial-article-files/46195}{Primary report}.
\bibitem{CGST} W.~Cook, A.~M.~H.~Gerards, A.~Schrijver and E.~Tardos,
\emph{Sensitivity theorems in integer linear programming}, Mathematical
Programming \textbf{34} (1986), 251--264.
\href{https://www.math.uwaterloo.ca/~bico/papers/sensitivity.pdf}{Author copy}.
\bibitem{THM} Y.~Takamatsu, S.~Hara and K.~Murota,
\emph{Continuous/Discrete Hybrid Convex Optimization and Its Optimality
Criterion}, Trans. Institute of Systems, Control and Information Engineers
\textbf{17} (2004), 409--411 (Japanese;
\href{https://kzmurota.fpark.tmu.ac.jp/paper/THM04hybridLenglish.pdf}{author English translation}).
\bibitem{Murota} K.~Murota, \emph{On steepest descent algorithms for discrete
convex functions}, SIAM J. Optimization \textbf{14} (2003), 699--707;
\href{https://www.keisu.t.u-tokyo.ac.jp/data/2002/METR02-12.pdf}{author report METR2002-12}, \S4.1.
\bibitem{Network} M.~Conforti, M.~Di~Summa, F.~Eisenbrand and L.~A.~Wolsey,
\emph{Network Formulations of Mixed-Integer Programs}, Mathematics of
Operations Research \textbf{34} (2009), 194--209;
\href{https://www.math.uwaterloo.ca/~bico/bellairs/network.pdf}{March 2008 author manuscript}.
\bibitem{Covers} M.~Conforti, L.~A.~Wolsey and G.~Zambelli,
\emph{Projecting an extended formulation for mixed integer covers on bipartite
graphs}, Mathematics of Operations Research (2010);
\href{https://personal.lse.ac.uk/Zambelli/papers/mix-tree.pdf}{author manuscript revised November 2009}.
\bibitem{Schrijver} A.~Schrijver, \emph{A combinatorial algorithm minimizing
submodular functions in strongly polynomial time}, J. Combinatorial Theory B
\textbf{80} (2000), 346--355.
\href{https://homepages.cwi.nl/~lex/files/minsubm6.pdf}{Author manuscript}, \S6.
\bibitem{GLS} M.~Gr\"otschel, L.~Lov\'asz and A.~Schrijver,
\emph{Geometric Algorithms and Combinatorial Optimization}, Springer, 1988.
\href{https://www.zib.de/userpage/groetschel/pubnew/paper/groetschellovaszschrijver1988.pdf}{Inspected 1988 edition}.
\end{thebibliography}
\end{document}
